% ============================================================
% f2014_hunan_science_q06_flowchart.tex
% 2014 湖南卷（理科）第 6 题：程序框图
% 来源 PDF：普通高考/2014/2014湖南理.pdf
% 原页：第 1 页（A3 横排）第 1 栏下部，印面 885，第 6 题题干下方
% 结构：开始 → 输入 t → 判断 t<0；
%       是：t=2t^2+1，回到判断框；否：S=t-3 → 输出 S → 结束。
% ============================================================
\begin{tikzpicture}[>=Stealth, line width=0.7pt, font=\small,
  x=1cm, y=1cm, line cap=round, line join=round]
  % 开始
  \draw[rounded corners=6pt] (-0.80,0.35) rectangle (0.80,-0.35);
  \node at (0,0) {开始};
  \draw[->] (0,-0.35) -- (0,-0.62);
  % 输入 t
  \draw (-1.15,-0.62) -- (1.15,-0.62) -- (1.00,-1.28) -- (-1.30,-1.28) -- cycle;
  \node at (0,-0.95) {输入 $t$};
  \draw[->] (0,-1.28) -- (0,-1.58);
  % 判断框
  \draw (0,-1.58) -- (1.95,-2.28) -- (0,-2.98) -- (-1.95,-2.28) -- cycle;
  \node at (0,-2.28) {$t<0$?};
  % 是：向左再向下进入循环体
  \draw[->] (-1.95,-2.28) -- (-2.85,-2.28) -- (-2.85,-3.25);
  \node at (-2.35,-2.12) {是};
  \draw (-3.80,-3.25) rectangle (-1.90,-3.85);
  \node at (-2.85,-3.55) {$t=2t^2+1$};
  % 回边：循环体左侧向上，水平进入判断框左上边
  \draw[->] (-3.80,-3.55) -- (-4.45,-3.55) -- (-4.45,-1.80) -- (-0.61,-1.80);
  % 否：向右再向下
  \draw[->] (1.95,-2.28) -- (3.00,-2.28) -- (3.00,-3.12);
  \node at (2.42,-2.12) {否};
  \draw (2.05,-3.12) rectangle (3.95,-3.72);
  \node at (3.00,-3.42) {$S=t-3$};
  % 输出 S
  \draw[->] (3.00,-3.72) -- (3.00,-4.25) -- (0,-4.25) -- (0,-4.45);
  \draw (-1.35,-4.45) -- (1.35,-4.45) -- (1.05,-5.15) -- (-1.65,-5.15) -- cycle;
  \node at (-0.30,-4.80) {输出 $S$};
  \draw[->] (0,-5.15) -- (0,-5.45);
  % 结束
  \draw[rounded corners=6pt] (-0.80,-5.45) rectangle (0.80,-6.15);
  \node at (0,-5.80) {结束};
\end{tikzpicture}
